Dans une expérience de Young, deux fentes séparées de \qty{\v{sep}}{mm} sont éclairées par une lumière de longueur d'onde \qty{\v{lam}}{nm}, et l'écran est à \qty{\v{L}}{m}. On recouvre une des fentes d'une mince lame de verre d'épaisseur \qty{\v{ep}}{\micro\metre} et d'indice de réfraction \num{\v{n}}. La figure d'interférence se déplace sur l'écran.

\begin{enumerate}[label=\alph*)]
    \item De combien le chemin optique de la lumière qui traverse la lame augmente-t-il par rapport au même trajet dans l'air ?
    \item De combien de franges la figure se déplace-t-elle ?
    \item Quelle distance sur l'écran correspond à ce déplacement ? Dans quel sens la frange centrale se déplace-t-elle ?
\end{enumerate}
